\section{Conventions and inputs from \texorpdfstring{\cite{I}}{[I]}}\label{dp:sec:conv}

This section fixes notation and recalls the inputs from~\cite{I} used throughout. Other results of~\cite{I} are cited where they are used.

\emph{Standing assumptions.} Throughout, $K$ is a field with a nontrivial nonarchimedean absolute value $\abs\cdot$, with completion $\Kh$, and $\varpi\in K$ is fixed with $0<\abs\varpi<1$. The powers of $\varpi$ have distinct absolute values, so $K$ is infinite. If $K$ is discretely valued, $\vartheta\in K$ is fixed with $\abs{K^\times}=\abs\vartheta^{\Z}$ and $\abs\vartheta<1$. Put $\mathfrak c_K=\abs\vartheta$ in that case, and $\mathfrak c_K=1$ otherwise.

The letters $P,P'$ denote source spaces and $W,Y,Z$ target spaces. All are normed $K$-spaces, not necessarily complete or ultrametric. A normed space is \emph{ultrametric} if $\norm{x+y}\le\max(\norm x,\norm y)$ for all its vectors $x,y$. The completion of $W$ is $\widehat W$. Products and finite direct sums carry the maximum norm. The open ball of radius $R$ around $x$ is $B(x,R)$. The \emph{$\ell^1$-sum} of normed spaces $Y_i$ is the space of the families $(y_i)_i$, $y_i\in Y_i$, with $\sum_i\norm{y_i}<\infty$, and this sum as norm~\cite[Definition~14.15]{I}. For a set $J$ let $\ell^1(J)=\ell^1(J,K)$ be the $\ell^1$-sum of copies of $K$ indexed by $J$, and let $\ell^1_{00}(J)$ be its subspace of finitely supported vectors.

\emph{Tensors.} For $n\ge1$, $P^{\otimes n}$ is the algebraic tensor power. For $t\in P^{\otimes n}$ put
\[
 \pi(t)=\inf\sum_l\prod_{i=1}^n\norm{h_{l,i}},\qquad
 \pi_{\max}(t)=\inf\max_l\prod_{i=1}^n\norm{h_{l,i}},
\]
both infima over the representations $t=\sum_lh_{l,1}\otimes\dots\otimes h_{l,n}$. Here $\pi$ is the projective seminorm and $\pi_{\max}$ is an ultrametric seminorm with $\pi_{\max}\le\pi$. For $n=0$ we read $P^{\otimes0}=K$, $\pi=\pi_{\max}=\abs\cdot$, $h^{\otimes0}=1$, and $\norm h^0=1$ also at $h=0$.

The space $\Gamma^n(P)\subseteq P^{\otimes n}$ is the span of the powers $h^{\otimes n}$, $h\in P$. As $K$ is infinite, it is the space of symmetric tensors~\cite[Section~L.3, opening paragraph]{I}. We write $\Gamma^n_\pi(P)$ for $\Gamma^n(P)$ with the restriction of $\pi$. On $\Gamma^n(P)$ we also use the seminorm
\begin{equation}\label{dp:eq:gamma-norm}
 \norm t_\Gamma=\inf\Bigl\{\sum_i\abs{c_i}\norm{h_i}^n:\ t=\sum_ic_i\,h_i^{\otimes n}\Bigr\},
\end{equation}
so that $\pi\le\norm\cdot_\Gamma$. If the bounded linear forms on $P$ separate points, then $\pi$ is a norm on $P^{\otimes n}$ and $\norm\cdot_\Gamma$ is a norm on $\Gamma^n(P)$. Indeed, the products of such forms are bounded for $\pi$, and they separate the points of $P^{\otimes n}$, as one checks on a finite-dimensional subspace containing the factors of a given tensor.

For $h_1,\dots,h_r\in P$ and integers $\alpha_i\ge0$ with $\abs\alpha=\alpha_1+\dots+\alpha_r=n$, the \emph{orbit sum} $O_\alpha(h_1,\dots,h_r)$ is the sum of the tensors $h_{i_1}\otimes\dots\otimes h_{i_n}$ over the tuples $(i_1,\dots,i_n)$ in which each $i$ occurs $\alpha_i$ times. With $\lambda^\alpha=\prod_i\lambda_i^{\alpha_i}$, expanding the power gives~\cite[(14.1)]{I}
\begin{equation}\label{dp:eq:orbit}
 \Bigl(\sum_{i=1}^r\lambda_ih_i\Bigr)^{\otimes n}=\sum_{\abs\alpha=n}\lambda^\alpha\,O_\alpha(h_1,\dots,h_r)\qquad(\lambda\in K^r),
\end{equation}
and $O_\alpha(h_1,\dots,h_r)\in\Gamma^n(P)$, since each coefficient of a polynomial with values in a subspace lies in that subspace, by interpolation.

For integers $n,n'\ge0$, $s\in P^{\otimes n}$ and $t\in P^{\otimes n'}$, the \emph{shuffle product} $s\cdot t\in P^{\otimes(n+n')}$ is the sum of the $\binom{n+n'}n$ tensors obtained from $s\otimes t$ by the permutations of the $n+n'$ factors that keep the order of the first $n$ factors and the order of the last $n'$. For $n=0$ or $n'=0$ the tensor $s\otimes t$ is the scalar multiple. The product is bilinear. Permutations of the factors preserve $\pi$, and $\pi(s\otimes t)\le\pi(s)\pi(t)$, so $\pi(s\cdot t)\le\binom{n+n'}n\pi(s)\pi(t)$. For $h,h'\in P$ we have $h^{\otimes n}\cdot h'^{\otimes n'}=O_{(n,n')}(h,h')$. So, by bilinearity and since orbit sums lie in $\Gamma^{n+n'}(P)$, the product maps $\Gamma^n(P)\times\Gamma^{n'}(P)$ into $\Gamma^{n+n'}(P)$. In the completed projective tensor powers it induces the shuffle product $*$ of~\cite[Section~L.3]{I}.

\emph{Enlarging the target.} Bounded polynomial laws, or \emph{laws}, their norm $\norm a$, maps ambient of degree $n$, and DP germs with radius, bound and coefficient laws are as in Definition~\ref{dp:def:laws}, which repeats~\cite[Definition~14.8]{I}. A law $a\colon P\to W$ of degree $n$ has a unique linear map $L_a\colon\Gamma^n(P)\to W$ with $L_a(h^{\otimes n})=a(h)$ for $h\in P$~\cite[Lemma~Q.1]{I}. Its \emph{ambient constant} is $\operatorname{amb}(a)=\sup\{\norm{L_a(t)}/\pi(t):t\in\Gamma^n(P)\}\in[0,\infty]$, where a quotient with denominator $0$ is read as $\infty$ if its numerator is positive and as $0$ if its numerator is $0$. The law $a$ is ambient of degree $n$ if and only if $\operatorname{amb}(a)<\infty$, and $\norm a\le\operatorname{amb}(a)$~\cite[text after Definition~14.8]{I}.

Let $U\subseteq P$ be open, $f\colon U\to W$, $x_0\in U$, $R>0$ with $B(x_0,R)\subseteq U$, and $M\ge0$. By~\cite[Definition~14.7]{I}, $f$ is \emph{ambient analytic}, or \emph{AA}, at $x_0$ with radius $R$ and bound $M$ if there are a linear isometry $j\colon W\to Z$ into a normed space and bounded $n$-linear maps $\beta_n\colon P^n\to Z$, $n\ge0$, with $\norm{\beta_n}\le MR^{-n}$, with $\beta_n(h,\dots,h)\in j(W)$ for all $h\in P$, and with $jf(x_0+h)=\sum_n\beta_n(h,\dots,h)$ for $\norm h<R$. It is AA at $x_0$ if this holds for some $R,M$, and AA on $U$ if it is AA at every point of $U$. For $Z=W$ and $j=\id_W$ we say that $f$ is \emph{analytic} at $x_0$ with radius $R$ and bound $M$, analytic at $x_0$ with radius $R$ if this holds for some $M$, and analytic at $x_0$ if it holds for some $R,M$. In~\cite{I} the bound $M$ and the maps $\beta_n$ are written $\mathfrak b$ and $b_n$.

By~\cite[Proposition~14.9]{I}: (1) $f$ is AA at $x_0$ with radius $R$ and bound $M$ if and only if it is DP at $x_0$ with radius $R$ and its coefficient laws $a_n$ satisfy $\operatorname{amb}(a_n)\le MR^{-n}$ for all $n$; it is then DP with bound $M$, and $ja_n(h)=\beta_n(h,\dots,h)$ for every $j$ and $(\beta_n)_n$ as above with this radius and any bound. (2) A law $a$ of degree $n$ such that $j\circ a$ is analytic at one point, for a linear isometry $j$ of $W$ into a normed space, is ambient of degree $n$.

Composition fails~\cite[Theorem~14.17]{I}. Let $\operatorname{char}K=p>0$, with $K$ complete or not, and let $f_X\colon H^{\rm at}\to W^{\rm at}$ and $g_X\colon W^{\rm at}\to Y^{\rm at}$ be the maps of~\cite[Definition~14.16]{I}, with $X=X^{\rm at}$ there. Then $f_X$ is AA at every point, $g_X$ is the diagonal of a bounded bilinear map, and $g_X\circ f_X$ is a law of degree $p+1$, DP everywhere and AA at no point. Whether composition fails in characteristic $0$ is open~\cite[Remark~14.20]{I}.

\emph{Tame pairs.} For $n\ge1$ let $\mathfrak a_n(P,W)$ be the supremum of $\operatorname{amb}(a)/\norm a$ over the nonzero laws $a\colon P\to W$ of degree $n$, and $0$ if there are none. The pair $(P,W)$ is \emph{tame} if $\sup_{n\ge1}\mathfrak a_n(P,W)^{1/n}<\infty$, and $P$ is \emph{tame} if $(P,W)$ is tame for every normed $W$. Put $\mathfrak a_n(P)=\sup\{\norm t_\Gamma/\pi(t):t\in\Gamma^n(P)\}$, with the same reading of quotients~\cite[Definition~14.36]{I}. Then $\mathfrak a_n(P,W)\le\mathfrak a_n(P)$, and $P$ is tame if and only if $\sup_{n\ge1}\mathfrak a_n(P)^{1/n}<\infty$~\cite[text after Definition~14.36]{I}.

For $W$ Banach, $(P,W)$ is tame if and only if every map from an open subset of $P$ into $W$ that is DP at a point is AA there~\cite[Proposition~14.37]{I}. By~\cite[Theorem~14.38]{I}: (1) if $\mathfrak a_n(P,W)\le\mathfrak q_0\mathfrak q^n$ for $n\ge1$, with $\mathfrak q_0,\mathfrak q\ge1$, then a germ from $P$ to $W$ that is DP at $x_0$ with radius $R$ and bound $M$ is AA there with radius $R/\mathfrak q$ and bound $\mathfrak q_0M$; (2) if $(P,W)$ is tame, $f$ from an open subset of $P$ into a normed space $Y$ is DP at $x_0$, and $g$ from an open subset of $Y$ into $W$ is DP at $f(x_0)$, then $g\circ f$ is AA at $x_0$; (3) if $P$ is Banach, $(P,P)$ is tame, and $f$ from an open subset of $P$ into $P$ is DP at $x_0$, with coefficient law $a_1$ bijective with bounded inverse, then $f$ has a local inverse germ, which is AA at $f(x_0)$ and at every point of a ball around it.

The examples are in~\cite[Theorem~14.39]{I}. (1) If $P,W$ are ultrametric, then $\mathfrak a_n(P,W)\le(\abs\varpi\mathfrak c_K)^{-n}$. (2) If $P=\ell^1(J)$ or $P=\ell^1_{00}(J)$, then $\mathfrak a_n(P)\le(8\mathrm e/\abs\varpi^2)^n$, with $\mathrm e$ Euler's number. (3) If $\operatorname{char}K=0$ and $P$ is ultrametric, then $\sup_{n\ge1}\mathfrak a_n(P)^{1/n}<\infty$, so $P$ is tame. (4) If $W$ is ultrametric and $P$ is any normed space, then every law from $P$ to $W$ is ambient, and $\mathfrak a_n(P,W)<\infty$ in each degree; whether $(P,W)$ is always tame is open. Between ultrametric spaces the AA germs and the DP germs coincide~\cite[Theorem~14.39(1)]{I}.

\emph{Alternating maps.} For normed spaces $E,D,F,F'$ and $k\ge1$, $\cL(E,D)$ is the normed space of bounded linear maps $E\to D$ with the operator norm, $\Mult^k(E;F)$ is the normed space of bounded $k$-linear maps $E^k\to F$, and $\Alt^k(E;F)$ is its subspace of alternating maps, those vanishing whenever two arguments agree. For $k=0$ both spaces are $F$. The inclusion $j_{E,F}\colon\Alt^k(E;F)\to\Mult^k(E;F)$ is a linear isometry with closed range~\cite[Section~2]{I}. The alternating degree $k$ is Bourbaki's $d$ in~\cite[n\textsuperscript{o}~7.8]{Bou}. Following~\cite[Section~2]{I}, for $u\in\cL(E,D)$, $v\in\cL(F,F')$ and $\beta\in\Mult^k(D;F)$ the \emph{joint action} at $z_1,\dots,z_k\in E$ is
\[
 B^k_{E,D;F,F'}(u,v)(\beta)(z_1,\dots,z_k)=v\bigl(\beta(uz_1,\dots,uz_k)\bigr),
\]
written $B$ in~\cite{I}. In the later sections $B(u,v)$, with operators as arguments, is this action, and $B(x,R)$, with a point and a radius, is the ball. The precomposition is the map $A^k_{E,D;F}$ from $\cL(E,D)$ to $\cL(\Alt^k(D;F),\Alt^k(E;F))$ that sends $u$ to the restriction of $B^k_{E,D;F,F}(u,\id_F)$ to alternating maps. By~\cite[Proposition~2.2]{I}, the joint action is the diagonal of a bounded $(k+1)$-linear map, and $A^k_{E,D;F}$ followed by $j_{E,F}\circ{-}$ is the diagonal of a bounded $k$-linear map of norm at most $1$. The projective exterior power $\Lambda^k_\pi E$ is the normed quotient of the algebraic exterior power by the null space of the seminorm built from wedges as $\pi$ is built from tensors, $\widehat\Lambda^k_\pi E$ is its completion, and $\Alt^k(E;F)\cong\cL(\Lambda^k_\pi E,F)$ isometrically, with precomposition by $u$ becoming precomposition by $\Lambda^k_\pi u$~\cite[Section~3.1, Proposition~3.1]{I}.

\emph{Bundles.} Let $\mathsf{Vec}_K$ be the category of normed $K$-spaces and bounded linear maps, with hom spaces $\cL(E,D)$. Let $m\in\N\cup\{\infty,\omega\}$; this is the index written $n$ in~\cite{I}. The class $C^m$ is as in the conventions of~\cite[Section~1]{I}, and $C^0$ means continuous. Fix a $C^m$ manifold $\mathcal M$ modeled on $P$. Then $\mathsf{VB}^m_K(\mathcal M)$ is the category of $C^m$ normed vector bundles over $\mathcal M$, with morphisms over $\id_{\mathcal M}$ whose local operator-valued expressions $x\mapsto T(x)\in\cL(E,F)$ are $C^m$. Let $\epsilon=(\epsilon_1,\dots,\epsilon_q)$ specify a variance in each variable. Put $\mathbf C^\epsilon=\mathsf C_1\times\dots\times\mathsf C_q$, with $\mathsf C_i=\mathsf{Vec}_K$ or $\mathsf{Vec}_K^{\mathrm{op}}$, and let $\mathbf B^\epsilon_{\mathcal M}$ be the corresponding product of bundle categories, with the same variances. A functor $\mathcal F\colon\mathbf C^\epsilon\to\mathsf{Vec}_K$ is $C^m$ if all its maps on products of hom spaces are jointly $C^m$. Let $\operatorname{Fun}^{C^m}(\mathbf C^\epsilon,\mathsf{Vec}_K)$ be the full subcategory of the functor category on these functors, and write $\eta$ for a natural transformation between two of them. The statement of~\cite[Theorem~2.1]{I} is this.
\begin{quote}
For a fixed $C^m$ manifold $\mathcal M$, fiberwise application defines, up to canonical natural isomorphism, a functor
\[
\begin{aligned}
 \mathcal B^m_{\mathcal M}\colon
 \operatorname{Fun}^{C^m}(\mathbf C^\epsilon,\mathsf{Vec}_K)
 &\longrightarrow
 \operatorname{Fun}(\mathbf B^\epsilon_{\mathcal M},\mathsf{VB}^m_K(\mathcal M)),\\
 \mathcal F&\longmapsto\mathcal F_{\mathcal M},\qquad
 \eta\longmapsto\eta_{\mathcal M}.
\end{aligned}
\]
The resulting bundles, bundle morphisms and natural transformations are given fiberwise by $\mathcal F$, its action on maps, and the components of $\eta$, respectively. The same conclusion holds if regularity on whole hom spaces is replaced by preservation of $C^m$ families parametrized by open subsets of the model space $P$.
\end{quote}
Preservation of $C^m$ families means that $\mathcal F\circ\gamma$ is $C^m$ for each $C^m$ map $\gamma$ from an open subset of $P$ into a product of hom spaces of $\mathbf C^\epsilon$. The theorem also holds for restricted categories of fibers containing the bundle morphisms~\cite[text after Theorem~2.1]{I}.

The proof in~\cite{I} uses of the class $C^m$ only these four properties. (i) Maps of the class on open sets compose, tuples of them into products are in the class, and the evaluation of operators is in the class. (ii) If $x\mapsto S(x)$ and $x\mapsto T(x)$ are local operator families of the class, so is $x\mapsto S(x)T(x)$. (iii) Constants are in the class. (iv) The class passes to restrictions to open subsets of $P$, and membership depends only on the map near each point. Step~1 there composes the transitions with the maps of the functor on hom spaces and evaluates operators, by (i), and glues chart by chart, by (iv). Step~2 applies the maps of the functor to local operator families, by (i), multiplies by operator families for changes of trivialization, by (ii), and works chart by chart, by (iv). Step~3 uses (iii), and Step~4 uses (iv). The category $\mathsf{VB}^m_K(\mathcal M)$ itself uses (ii) for the composition of morphisms and (iii) for identities. Corollary~\ref{dp:cor:bundles} applies this to DP maps.
